% test-018.tex -- sesion unica para verificar silent-resto.patch
%
% Este archivo REEMPLAZA a las listas de casos de los tests 008, 010, 012
% y 017: un solo PDF, una sola grabacion.
%
% Orden de prioridad:
%   A (casos 1 a 11)  : lo que decide si el parche funciono.
%   B (12 a 25)       : silent; no debe leerse NADA despues de la etiqueta.
%   C a G (26 a 59)   : rutas de codigo cuyo MathML cambio; deben leerse
%                       como esta anotado en "doc:" (o como antes del parche).
% Si hay que parar a la mitad, con A y B ya esta lo importante.
%
% Compilar:  lualatex-dev test-018.tex
\DocumentMetadata
  {
    lang=es-CL,
    pdfversion=2.0,
    pdfstandard=ua-2,
    tagging=on,
    tagging-setup={math/setup={mathml-SE}},
  }
\documentclass{article}
\usepackage[spanish]{babel}
\usepackage{unicode-math}
\usepackage{spintent}
\usepackage{hyperref}
\hypersetup
  {
    pdfdisplaydoctitle = true,
    pdftitle           = {test-018: sesion unica de silent-resto},
  }
\begin{document}

\newcommand{\xcuad}{x^2}%
\spnewread{\cong}{es-congruente-con}%
\spnewread{\perp}{es-perpendicular-a}%
\spnewread{\approx}{es-aproximadamente-igual-a}%
\spnewread{\neq}{es distinto de}%
\spnewread{\geq}{es_mayor_o_igual_que}%
\spnewread{\in}{pertenece a}%
\spnewread{\xcuad}{x al cuadrado}%


\section*{A. Lo que decide el parche}

% doc: «»
1. Silent con resultado, divisores de 6: $\spnD[silent=true,result=true](6)$.

% doc: «»
2. Silent con resultado y result-num, divisores de 12: $\spnD[silent=true,result=true,result-num=3](12)$.

% doc: «»
3. Silent con resultado, múltiplos de 6: $\spnM[silent=true,result=true](6)$.

% doc: «»
4. Silent con resultado y result-num, múltiplos de 12: $\spnM[silent=true,result=true,result-num=3](12)$.

% doc: «»
5. Silent con resultado, máximo común divisor: $\spmcd[silent=true,result=true](12, 18)$.

% doc: «»
6. Silent con resultado, mínimo común múltiplo: $\spmcm[silent=true,result=true](12, 18)$.

% doc: «divisores de 6 es un subconjunto de divisores de 12»
7. Divisores incluidos: $\spnD(6) \subset \spnD(12)$.

% doc: «divisores de 6 son 1 2 3 y 6 es un subconjunto de divisores de 12 son 1 2 3 4 6 y 12»
8. Divisores con resultado, incluidos: $\spnD[result=true](6) \subset \spnD[result=true](12)$.

% doc: «múltiplos de 4 son 4 8 12 puntos suspensivos intersección múltiplos de 6 son 6 12 18 puntos suspensivos»
9. Múltiplos con resultado, intersección: $\spnM[result=true,result-num=3](4) \cap \spnM[result=true,result-num=3](6)$.

% doc: «punto A es igual a 1 coma 2»
10. Punto igual a coordenadas: $\spPoint{A} = \spcoord{1, 2}$.

% doc: «vector a es igual a 1 coma 2»
11. Tres vectores con coordenadas, cada uno en su fórmula: $\spvec{a}(1,2)$, $\spvec{b}(3,4)$ y $\spvec{c}(4,6)$.

\section*{B. Con silent true no debe leerse nada}

% doc: «»
12. Silent en sprecta con dos puntos con subíndices: $\sprecta[silent=true]{A_{1}B_{2}}$.

% doc: «»
13. Silent en sprecta con una letra con subíndice: $\sprecta[silent=true]{r_{1}}$.

% doc: «»
14. Silent en spcirc con un centro con subíndice: $\spcirc[silent=true](O_{1})$.

% doc: «»
15. Silent en spcirc con radio entero: $\spcirc[silent=true](O,3)$.

% doc: «»
16. Silent en spcirc con radio segmento: $\spcirc[silent=true](O,AB)$.

% doc: «»
17. Silent en spAcirc con un centro con subíndice: $\spAcirc[silent=true](O_{1})$.

% doc: «»
18. Silent en spAcirc con un centro con prima: $\spAcirc[silent=true](O')$.

% doc: «»
19. Silent en spPcirc con un centro con subíndice: $\spPcirc[silent=true](O_{1})$.

% doc: «»
20. Silent en spnD con un número: $\spnD[silent=true](6)$.

% doc: «»
21. Silent en spnM con un número: $\spnM[silent=true](12)$.

% doc: «»
22. Silent en spmcd con dos números: $\spmcd[silent=true](12, 18)$.

% doc: «»
23. Silent en spmcm con dos letras: $\spmcm[silent=true](m, n)$.

% doc: «»
24. Silent en spread con un comando propio: $\spread[silent=true]{\xcuad}$.

% doc: «»
25. Silent en spread con un símbolo: $\spread[silent=true]{\cong}$.

\section*{C. spnD y spnM sin silent}

% doc: «divisores de 6»
26. Divisores de un número: $\spnD(6)$.

% doc: «divisores de 12»
27. Divisores de un número de dos cifras: $\spnD(12)$.

% doc: «divisores de n»
28. Divisores de una letra: $\spnD(n)$.

% doc: «múltiplos de 6»
29. Múltiplos de un número: $\spnM(6)$.

% doc: «divisores de n»
30. Divisores sin argumento: $\spnD$.

% doc: «divisores de 12 son 1 2 3 4 6 y 12»
31. Divisores con resultado: $\spnD[result=true](12)$.

% doc: «divisores de 7 son 1 y 7»
32. Divisores con resultado de dos elementos: $\spnD[result=true](7)$.

% doc: «divisores de m es un subconjunto de divisores de n»
33. Divisores de dos letras en una fórmula: $\spnD(m) \subset \spnD(n)$.

\section*{D. spmcd y spmcm sin silent}

% doc: «máximo común divisor entre 12 y 18»
34. Máximo común divisor: $\spmcd(12, 18)$.

% doc: «mínimo común múltiplo entre 12 y 18»
35. Mínimo común múltiplo: $\spmcm(12, 18)$.

% doc: «máximo común divisor entre 12 y 18 es igual a 6»
36. Máximo común divisor con resultado: $\spmcd[result=true](12, 18)$.

% doc: «mínimo común múltiplo entre 12 y 18 es igual a 36»
37. Mínimo común múltiplo con resultado: $\spmcm[result=true](12, 18)$.

% doc: «máximo común divisor entre m y n por mínimo común múltiplo entre m y n es igual a m por n»
38. Producto de los dos con letras: $\spmcd(m, n) \spmsym \spmcm(m, n) = m \spmsym n$.

% doc: «máximo común divisor entre 12 y 18 es igual a 6 más mínimo común múltiplo entre 12 y 18 es igual a 36»
39. Suma de los dos con resultado: $\spmcd[result=true](12, 18) + \spmcm[result=true](12, 18)$.

\section*{E. spcoord y coordenadas}

% doc: «menos 3 coma 4»
40. Coordenadas con signo: $\spcoord{-3, 4}$.

% doc: «x subíndice 1 coma y subíndice 1»
41. Coordenadas con subíndices: $\spcoord{x_{1},y_{1}}$.

% doc: «3 coma 5 punto y coma 4 coma 2»
42. Coordenadas con coma decimal: $\spcoord{3,5; 4,2}$.

% doc: «el raíz cuadrada de 2 coma 1»
43. Coordenadas con una raíz: $\spcoord{\sqrt{2}, 1}$.

% doc: «punto A menos 3 coma 4»
44. Un punto con coordenadas: $\spPoint{A}(-3,4)$.

% doc: «punto A sub uno menos 3 coma 4»
45. Un punto con subíndice y coordenadas: $\spPoint{A_{1}}(-3,4)$.

% doc: «vector a es igual a menos 3 coma 4»
46. Un vector con coordenadas: $\spvec{a}(-3,4)$.

% doc: «vector A B es igual a punto B 1 coma 2 menos punto A 4 coma menos 2»
47. Un vector con coordenadas en una fórmula destacada:
\[ \spvec{AB} = \spPoint{B}(1,2) - \spPoint{A}(4,-2) \].

\section*{F. spread: el símbolo ahora va dentro de un contenedor}

% doc: «triángulo blanco apuntando hacia arriba a b c es congruente con triángulo blanco apuntando hacia arriba d e f»
48. Símbolo registrado en una fórmula: $\triangle ABC \spread{\cong} \triangle DEF$.

% doc: «a es distinto de b»
49. Lectura con espacios: $a \spread{\neq} b$.

% doc: «a es mayor o igual que b»
50. Lectura con guiones bajos: $a \spread{\geq} b$.

% doc: «x pertenece a a»
51. Otra lectura con espacios: $x \spread{\in} A$.

% doc: «a es perpendicular a b es aproximadamente igual a c»
52. Dos símbolos en una fórmula: $a \spread{\perp} b \spread{\approx} c$.

% doc: «x al cuadrado»
53. Un comando propio: $\spread{\xcuad}$.

% doc: «y es igual a x al cuadrado más 1»
54. El comando propio en una fórmula: $y = \spread{\xcuad} + 1$.

% doc: «a es paralela a b»
55. Un símbolo sin registrar con read-arg: $a \spread[read-arg={es paralela a}]{\parallel} b$.

\section*{G. Controles que no deben haber cambiado}

% doc: «circunferencia de centro en O y radio 3»
56. Control de spcirc con radio entero: $\spcirc(O,3)$.

% doc: «circunferencia de centro en O y radio medida del lado A B»
57. Control de spcirc con radio segmento: $\spcirc(O,AB)$.

% doc: «área de la circunferencia con centro en O sub uno»
58. Control de spAcirc con un centro con subíndice: $\spAcirc(O_{1})$.

% doc: «recta A sub uno B sub dos»
59. Control de sprecta con dos puntos con subíndices: $\sprecta{A_{1}B_{2}}$.

\end{document}
